
\documentclass{article}
\usepackage{graphicx}
\usepackage{amsmath}
\usepackage{amssymb}
\usepackage{amsfonts}
\usepackage{hyperref}
\usepackage{exercises}
\usepackage{indentfirst}
\usepackage{circledsteps}
\usepackage{amsthm}
\usepackage{csquotes}
\usepackage{yhmath}
\usepackage{comment}
\usepackage{placeins}
\usepackage{geometry}
\usepackage{cleveref}
\usepackage{stmaryrd}
\usepackage{hyperref}
\usepackage{tikz}
\usepackage{tabto}
\newcommand*{\oldneg}{\mathord{\sim}}
\newcommand{\Z}{\mathbb{Z}}
\newcommand{\N}{\mathbb{N}}
\let\oldemptyset\emptyset
\newcommand{\lbc}{\llbracket}
\newcommand{\rbc}{\rrbracket}
\newcommand{\ltop}[1]{\stackrel{\leftarrow}{#1}}

\theoremstyle{definition}
\newtheorem{definition}{Definition}
\newtheorem{theorem}{Theorem}[section]
\newtheorem{corollary}{Corollary}[theorem]
\newtheorem{lemma}[theorem]{Lemma}

\setlength{\parindent}{.5in}
\geometry{margin=0.5in}
\setlength{\parskip}{12pt}
\title{Chapter 19 Notes}
\begin{document}
\date{}
\maketitle

\section*{Well Ordering of a Set Implies Linear Ordering of the Power Set}
Suppose some set $S$ can be well ordered under some relation $\le$. Define the following ordering $\preceq$ on 
$\mathcal{P}(S)$: for any $x,y\in\mathcal{P}(S)$, $x\preceq y$ iff 
$x=y$ or the least element (under $\le$) $e$ of $x\triangle y$ (the symmtetric difference of $x$ and $y$) 
is in $x$ but not in $y$. We show that this is a linear ordering.

Reflexivity: This is obvious.

Anti-Symmetry: Suppose $x\preceq y$ and $y\preceq x$. Using $x\preceq y$, there are then two cases.

Case 1) $x=y$. This is just what we wanted to obtain.

Case 2) $e\in x$ and $e\notin y$ for $e$ the least element of $x\triangle y$. We also have $y\preceq x$. There are then two 
subcases based on this. The first is when $y=x$, and again, this is just the result we wanted. Otherwise, 
$e\in y$ and $e\notin x$ for $e$ the least element of $y\triangle x$. But this is a contradiction, as 
$e\in x$. Hence it must be that $x=y$.

Transitivity: Suppose $x\preceq y$ and $y\preceq z$. Using $x\preceq y$, there are then two cases.

Case 1) $x=y$. This obviously implies $x\preceq z$.

Case 2) $e\in x$ and $e\notin x$ for $e$ the least element of $x\triangle y$. There are then two subcases, based off of 
$y\preceq z$. In the first subcase, 
$y=z$. This obviously implies $x\preceq z$. In the second subcase, $f\in y$ and $f\notin z$ for $f$ the least element of 
$y\triangle z$. We aim to show that $h\in x$ and $h\notin z$ for $h$ the least element of $x\triangle z$. First, for any element 
$d<e$, if $d\in z$, then either $d\in y$ (which implies $d\in x$), or $f<d<e$, which would imply that $f\in x$ (and 
recall that $f\notin z$). What this says is that no 
element $d<e$ of $z$ could be the least element of $x\triangle z$. Then, if 
$e\notin z$, we immediately know that $e$ is the least element of $x\triangle z$. On the other hand, suppose that $e\in z$. 
This means, because $e\notin y$, that $f<e$. Then it would follow that $f\in x$. In this case, $f$ would be the least element of 
$x\triangle z$, as we have already shown there are no elements $d<e$ of $z$ in $x\triangle z$. Thus we have $x\preceq z$. 

Linearity: This is obvious.

\section*{$2^{|S|}=|\mathcal{P}(S)|$}
The result is easy to prove by induction if the cardinality of $S$ is finite; so suppose that the cardinality of $S$ is 
infinite. We show that, if $\mathcal{P}(S)$ has cardinality, then $2^{|S|}=|\mathcal{P}(S)|$, by showing that 
both $2^S\preceq\mathcal{P}(S)$ and $\mathcal{P}(S)\preceq 2^S$ (which works, as $2^{|S|}$ is equinumorous with $2^S$).

First we show that $2^S\preceq\mathcal{P}(S)$. 
Consider any function $f\in 2^S$. There are three cases: (1) either 
$f(x)=1$ for all $x\in\textrm{Dom}(f)$; (2) $f(x)=0$ for all $x\in\textrm{Dom}(f)$; or 
$f(x)=1$ for some $x\in\textrm{Dom}(f)$ but $f(x)=0$ for some other $x\in\textrm{Dom}(f)$. If we can show that the 
set of all functions of type (1) has cardinality $|\mathcal{P}(S)|$, and show the same result for the set of all 
functions of types (2) and (3), then we will have the desired result. It is obvious that there are 
$|\mathcal{P}(S)|$ functions of types (1) and (2). Now, consider any function $f$ of type (3). There are 
$|\mathcal{P}(S)|$ possible ways the subset of $\textrm{Dom}(f)$, where each $x$ in the subset is such that 
$f(x)=0$, could be made up. It suffices to show that, for each such possible subset, there are only 
$|\mathcal{P}(S)|$ possible ways the subset of $\textrm{Dom}(f)$, where each $x$ in the subset is such that 
$f(x)=1$, could be made up (because this would mean that the set of all functions of 
type (3) would have cardinality equal to $|\mathcal{P}(S)|\times|\mathcal{P}(S)|$, and this equals 
$|\mathcal{P}(S)|$ as $|\mathcal{P}(S)|$ is infinte, as $S$ was assumed to be infinite). But this is obvious. Hence 
we know that the set of all functions of type (3) is of size $|\mathcal{P}(S)|$, which completes this case. 

It is obvious that $\mathcal{P}(S)\preceq 2^S$, so this completes the proof.

\end{document}
